\scanpage{106}

et considérons le sous-schéma de $\check{\mathbb V}(\det\check V)=\mathbb V(\det(\check{\unsure{W|V}}))=\mathbb V(\det(W)^{-1})$ défini par l'équation

\[
\omega^{\otimes(-2)}=\delta'(q)
\]

C'est un torseur sous $\mu_{2\,S}\simeq(\mathbb Z/2)_S$, i.e. un rev. \uline{étale} de degré 2 de $S$, soit $\Omega$.

% slanted

\margin{car. rés. $\ne2$ !!}

[Dans le cas où $q$ est spéciale orthogonale, i.e. munie de $\omega_0$ telle que $\omega_0^{\otimes-2}=\underline{-}\delta'(q)$ -- attention au signe ! -- les $\omega$ correspondent aux $\lambda\omega_0$ avec $\lambda^2=-1$, donc $\Omega=S[\underline i]\overset{\text{déf}}{=}\operatorname{Spec}\mathcal O_S[T]/(T^2+1)$]

Considérons la clôture projective $\hat\Sigma$ $(=Q)$ de $\Sigma$, et \deleted{\unsure{$\Sigma$}} son lieu à l'infini $X^{\subset}$, qui est aussi la \deleted{\unsure{localité}}\inserted{conique} projective \inserted{$\subset H=\check{\mathbb P}(W)$, hyperplan à l'infini de $\hat E=\check{\mathbb P}(\unsure{\mathscr E|\mathcal E})$,} définie par la forme $q$.

On voit que, \deleted{\unsure{$\Sigma$}} si $\sigma$ désigne l'antipodisme (symétrie p.r. à $o$ dans $E$, \unsure{de même|ou même}, dans $\hat E$) on

\[
\bigl(\qquad\hat E^\sigma=\inserted{=o(S)\amalg H,\ \unsure{\therefore}} H=P(V)\text{ (hyperplan à l'infini de }\hat E)
\]

donc $\hat\Sigma^\sigma=H\cap\hat\Sigma=X$.

% archive sheet number: 105
